\documentclass[11pt]{article}
\usepackage[a4paper,margin=1in]{geometry}
\usepackage{amsmath,amssymb,amsthm}
\usepackage[T1]{fontenc}
\usepackage{lmodern}
\usepackage[hidelinks]{hyperref}
\setlength{\emergencystretch}{3em}

\theoremstyle{plain}
\newtheorem{theorem}{Theorem}
\newtheorem{lemma}[theorem]{Lemma}
\newtheorem{proposition}[theorem]{Proposition}
\newtheorem{corollary}[theorem]{Corollary}
\theoremstyle{definition}
\newtheorem{definition}[theorem]{Definition}
\theoremstyle{remark}
\newtheorem{remark}[theorem]{Remark}

\newcommand{\F}{\mathbb F}

\title{A Counterexample to Conjecture 00000001063:\\
Kakeya Sets in $\F_3^2$ Have at Least Seven Points}
\author{%
  Feiyu\thanks{Independent researcher.}
  \and
  Claude (Opus 5.5)\thanks{Anthropic.}%
}
\date{2026-10-02}

\begin{document}
\maketitle

\begin{abstract}
Conjecture 00000001063 asserts that the minimal size of a Kakeya set in the
plane $\F_q^2$ is exactly $q(q+1)/2$. This fails for $q=3$: the conjecture
predicts $6$, but every Kakeya set in $\F_3^2$ has at least $7$ points, and $7$
is attained. The refutation is formalized in Lean~4 against Mathlib.
\end{abstract}

\section{The conjecture as stated}

The original text of conjecture 00000001063 reads:

\begin{quote}
Definition: A Kakeya set in $\F_q^2$ contains a line in every direction.
Conjecture: The minimal size of a Kakeya set is exactly $q(q+1)/2$, attained by
an affine-translation construction taking one segment on each of $q$ lines of
increasing slope; that is, the optimal constant in Dvir's lower bound $c\cdot
q^2$ is $c = 1/2$.
\end{quote}

\section{Reading of the statement}

Let $\F$ be a finite field with $q$ elements. A \emph{line} of $\F^2$ in the
direction $v\in\F^2\setminus\{0\}$ is a set $\{a+tv : t\in\F\}$ with
$a\in\F^2$.

\begin{definition}
A set $K\subseteq\F^2$ is a \emph{Kakeya set} if for every $v\ne0$ there is an
$a\in\F^2$ with $a+tv\in K$ for all $t\in\F$.
\end{definition}

This is the standard notion in the finite field Kakeya problem
\cite{dvir,bm}, and it is the one in the conjecture's own definition. The size
of $K$ is its number of points. The statement places no restriction on $q$, so
its first sentence asserts, for every finite field $\F_q$:
\begin{equation}\label{eq:claim}
  \text{some Kakeya set in }\F_q^2\text{ has exactly }q(q+1)/2\text{ points,
  and none has fewer.}
\end{equation}
We show that \eqref{eq:claim} is false for $q=3$. The remaining clauses of the
conjecture (how the minimum is attained, and the asymptotic constant $c=1/2$)
are offered as consequences of \eqref{eq:claim}; the refutation does not use
them, and we comment on them in Section~\ref{sec:remarks}.

\section{Result}

\begin{theorem}\label{thm:main}
Every Kakeya set in $\F_3^2$ has at least $7$ points, and there is a Kakeya set
in $\F_3^2$ with exactly $7$ points. Consequently the minimal size of a Kakeya
set in $\F_3^2$ is $7\ne 6=3\cdot4/2$, and Conjecture 00000001063 is false.
\end{theorem}

\section{Proof}

Throughout, $\F=\F_3=\{0,1,2\}$, in which $3=0$, $-2=1$ and $2\cdot2=1$.

\paragraph{Reduction to four lines.}
Up to nonzero scalars there are four directions,
$(1,0)$, $(0,1)$, $(1,1)$ and $(1,2)$. A line in each of them is cut out by an
equation of the form
\[
  y=c_1,\qquad x=c_2,\qquad y=x+c_3,\qquad y=2x+c_4
\]
respectively: the line $\{a+t(1,0)\}$ is $y=a_2$, the line $\{a+t(0,1)\}$ is
$x=a_1$, the line $\{a+t(1,1)\}$ is $y=x+(a_2-a_1)$, and the line
$\{a+t(1,2)\}$ is $y=2x+(a_2-2a_1)$. So a Kakeya set $K$ contains, for some
$c_1,\dots,c_4\in\F_3$, the union
\[
  U(c_1,c_2,c_3,c_4)=\{(x,y) : y=c_1\ \text{or}\ x=c_2\ \text{or}\ y=x+c_3\
  \text{or}\ y=2x+c_4\}.
\]

\begin{lemma}\label{lem:seven}
$|U(c_1,c_2,c_3,c_4)|\ge7$ for all $c_1,c_2,c_3,c_4\in\F_3$.
\end{lemma}

\begin{proof}
Call the four lines $L_1,\dots,L_4$. Lines with different directions meet in
exactly one point, so each of the six pairs $L_i\cap L_j$ is a single point and
inclusion--exclusion gives $|U|\ge 4\cdot3-6=6$, with equality only if the six
intersection points are pairwise distinct, that is, only if no point lies on
three of the lines. Suppose that is the case. Then on each line $L_i$ the three
points $L_i\cap L_j$ ($j\ne i$) are distinct, so they are all three points of
$L_i$.

On $L_1$ ($y=c_1$) these points have $x$-coordinates $c_2$, $c_1-c_3$ and
$2(c_1-c_4)$; being all of $\F_3$, they sum to $0$:
$c_2+c_1-c_3+2c_1-2c_4=c_2-c_3+c_4=0$.
On $L_2$ ($x=c_2$) the points have $y$-coordinates $c_1$, $c_2+c_3$ and
$2c_2+c_4$, which likewise sum to $0$: $c_1+c_3+c_4=0$.
Subtracting, $c_2-c_1-2c_3=c_2-c_1+c_3=0$, that is $c_1-c_3=c_2$. But $c_1-c_3$
and $c_2$ are the $x$-coordinates of $L_1\cap L_3$ and $L_1\cap L_2$, which
both lie on $L_1$; so these two points coincide and lie on three lines, a
contradiction. Hence $|U|\ge7$.
\end{proof}

The Lean development does not rely on this argument: it checks
Lemma~\ref{lem:seven} by evaluating $|U(c_1,c_2,c_3,c_4)|$ for all $3^4=81$
choices in the kernel.

\begin{proof}[Proof of Theorem~\ref{thm:main}]
If $K$ is a Kakeya set, it contains some $U(c_1,c_2,c_3,c_4)$, so $|K|\ge7$ by
Lemma~\ref{lem:seven}. Conversely, the $7$-point set
\[
  K_7=\{(0,0),(1,0),(2,0),(0,1),(0,2),(1,1),(2,2)\}
\]
is a Kakeya set: it contains the lines $y=0$, $x=0$ and $y=x$ through the
origin, and the line $y=2x+1=\{(0,1),(1,0),(2,2)\}$. Hence the minimal size is
$7$. The conjecture predicts $q(q+1)/2=6$ for $q=3$, and no Kakeya set of that
size exists.
\end{proof}

\section{Formalization}

The accompanying Lean~4 project (\texttt{lean/Submission/Basic.lean}) states
the conjecture and proves its negation against Mathlib.

\begin{itemize}
  \item \texttt{IsKakeya K} --- for every nonzero $v\in F\times F$ there is $a$
    with $a+t\bullet v\in K$ for all $t\in F$, over an arbitrary field $F$.
  \item \texttt{MinKakeyaSizeIs F s} --- some Kakeya set in $F\times F$ has
    exactly $s$ points and none has fewer.
  \item \texttt{ConjectureHolds} --- for every finite field $F$,
    \texttt{MinKakeyaSizeIs F (q * (q + 1) / 2)} with
    $q=\texttt{Fintype.card F}$; this is \eqref{eq:claim}.
  \item \texttt{fourLines c1 c2 c3 c4} --- the set $U(c_1,c_2,c_3,c_4)$ over
    \texttt{ZMod 3}.
  \item \texttt{seven\_le\_card\_fourLines} --- Lemma~\ref{lem:seven},
    proved by \texttt{decide}.
  \item \texttt{exists\_fourLines\_subset} --- a Kakeya set in
    $\mathbb Z/3\times\mathbb Z/3$ contains some $U(c_1,\dots,c_4)$;
    \texttt{seven\_le\_card\_of\_isKakeya} --- hence it has at least $7$ points.
  \item \texttt{K7}, \texttt{isKakeya\_K7}, \texttt{card\_K7},
    \texttt{minKakeyaSizeIs\_zmod\_three} --- the minimal size over $\F_3$ is
    exactly $7$.
  \item \texttt{conjecture\_00000001063\_false} ---
    $\lnot\,$\texttt{ConjectureHolds}, by specializing to
    $F=\texttt{ZMod 3}$.
\end{itemize}

The project builds with \texttt{lake build} on
\texttt{leanprover/lean4:v4.33.1} with Mathlib \texttt{v4.33.1}. It contains
no \texttt{sorry}, no \texttt{admit} and no \texttt{native\_decide}, and
introduces no axioms: \texttt{\#print axioms
conjecture\_00000001063\_false} reports exactly \texttt{propext},
\texttt{Classical.choice} and \texttt{Quot.sound}.

\section{Remarks}\label{sec:remarks}

The value $7$ found above is the case $q=3$ of a theorem of Blokhuis and
Mazzocca \cite{bm}: for odd $q$ the minimal size of a Kakeya set in $\F_q^2$ is
\[
  \frac{q(q+1)}2+\frac{q-1}2,
\]
while for even $q$ it is $q(q+1)/2$. So the exact formula of the conjecture is
correct for even $q$ and wrong for every odd $q$; for instance a brute-force
search gives $17$, not $15$, for $q=5$. The bound $q(q+1)/2$ itself is the
easy lower bound obtained, as in Lemma~\ref{lem:seven}, from inclusion--exclusion
over $q+1$ lines in distinct directions, and it is attained exactly when no
three of the lines are concurrent, which is possible only for even $q$.

The asymptotic clause of the conjecture is consistent with these values: the
minimal size is $q^2/2+O(q)$ for all $q$, so the optimal constant in a lower
bound of the form $c\,q^2$ is indeed $1/2$. It is the exact statement
\eqref{eq:claim} that fails. Conjecture 00000001064, which immediately follows
this one, states the correct value for odd $q$.

\begin{thebibliography}{9}
\bibitem{bm} A.~Blokhuis and F.~Mazzocca, \emph{The finite field Kakeya
  problem}, in: Building Bridges, Bolyai Society Mathematical Studies~19,
  Springer, 2008, pp.~205--218. arXiv:0911.4370.
\bibitem{dvir} Z.~Dvir, \emph{On the size of Kakeya sets in finite fields},
  J. Amer. Math. Soc. 22 (2009), 1093--1097.
\end{thebibliography}

\end{document}
